% ===================================================================== Chapter 1
\chapter{Introduction}\label{ch:intro}

\section{Background}
The internship project \emph{Flow Behavior and Thermal Effects in Multiphysics Systems} was carried out at Eleation between February and
May 2025 with ANSYS Workbench, Fluent and Mechanical. Its original files were not retained, so the analysis reported here was carried
out again after the internship: a representative multiphysics problem was defined, solved and audited from first principles, with every
input labelled as selected or assumed and every result traceable to a calculation or simulation in the project record. The note on the
analysis record at the front of the report states this formally; it applies to every number in the following chapters.

\section{Multiphysics Engineering Problems}
Many aerospace and energy components are loaded by several interacting physical processes at once. A heated duct, a combustor liner or
a heat-exchanger tube carries a flow that sets the convective boundary condition, the convective boundary condition sets the wall
temperature, and the wall temperature sets thermal strain and, where that strain is prevented, thermal stress. Each link has its own
governing equations and numerical methods, and the uncertainty of the final structural quantity depends on every link in the chain.

\section{Importance of Coupled Flow and Thermal Analysis}
In a conjugate problem the wall temperature is part of the solution rather than a boundary condition. Its accuracy depends on how well
the near-wall flow is resolved, because the air-side film usually carries most of the thermal resistance. A prescribed wall temperature
or an assumed heat-transfer coefficient would decouple the structural result from the flow physics that actually determines it.
Conjugate heat transfer (CHT) with a wall-resolved turbulence model therefore forms the thermal basis of this study.

\section{Motivation}
The study has two purposes. The first is technical: to demonstrate a complete and auditable simulation chain from flow to
structural stability, including verification at each step (analytical baseline, conservation checks, mesh studies, mapping checks).
The second is documentary: to provide a reproducible and auditable record of the internship problem, in which the status of every
value is stated explicitly.

\section{Problem Statement}
An air-cooled Inconel 718 duct is heated uniformly on its outer surface. The central question is:
\begin{quote}
\emph{Does thermal stress in the heated duct come mainly from the temperature gradient through the wall, or from restraint of the
duct's overall thermal growth -- and what are the consequences of that restraint for stability?}
\end{quote}
Answering it requires the flow field and heat-transfer coefficient, the conjugate temperature field in the wall, the thermal stresses for
free and restrained expansion, and the linear stability of the restrained duct, together with their sensitivity to the operating
conditions, the wall thickness and the end supports.

\section{Scope}
The study covers a steady-state, one-way coupled chain: Fluent CFD with conjugate heat transfer, transfer of the converged solid
temperature field to Mechanical, linear static analysis of two load cases (LC1 free expansion, LC2 axially restrained), linear
eigenvalue buckling on the LC2 pre-stress, a one-factor-at-a-time parametric study and a support-sensitivity study. Outside the scope are:
\begin{itemize}
\item transient start-up and thermal cycling;
\item radiation;
\item two-way fluid--structure coupling;
\item geometric or material nonlinearity and imperfection-sensitive buckling;
\item experimental validation.
\end{itemize}
These exclusions are revisited as limitations in Chapter~\ref{ch:unc} and as future work in Chapter~\ref{ch:future}. An additional
geometrically nonlinear, elastic buckling analysis with ANSYS LS-DYNA, performed as an extension of the thermo-structural analysis, is
reported separately in Section~\ref{sec:lsdyna}; it does not change the scope or the results of the main study.

\section{Objectives}
\begin{enumerate}
\item Define a representative heated-duct problem and an independent analytical baseline before any simulation.
\item Solve the turbulent flow and conjugate heat transfer in ANSYS Fluent to a converged, conservative solution with a resolved wall.
\item Quantify the CFD discretisation error with a three-mesh study and select the mesh for the structural phase.
\item Transfer the solid temperature field to ANSYS Mechanical and verify the transfer.
\item Compute thermal stress and deformation for free and restrained expansion and check structural mesh sensitivity.
\item Evaluate the linear buckling behaviour of the restrained duct and its dependence on the end supports.
\item Quantify the effect of inlet velocity, heat flux and wall thickness with actual CFD, structural and buckling solutions.
\item Record every uncertainty separately and state which conclusions the evidence supports.
\end{enumerate}
